\section{Proof of Theorem~\ref{thm:half}}\label{sec:halfproof}

Let $D_{*}$ and $C_0$ be as in Lemma~\ref{lem:affine}, and let $K_1$ be as in
Theorem~\ref{thm:raise}. Let $G$ and $P$ be as in Theorem~\ref{thm:half}, and suppose that
$q(G,P)>\half$. We derive a contradiction when $r$ is large. All $o(\cdot)$ refer to $r\to\infty$ and
are uniform in $G$ and $P$. Put
\[
  x=\log r,\qquad z=\log x,\qquad A=6\,e^{168},\qquad L=\lfloor Ax\rfloor,\qquad
  t=\lfloor z-2\log z\rfloor .
\]

\paragraph{Two girth scales.}
Let $D_j=2^{j}D_{*}$ and let $S_0'\subseteq S_1'\subseteq\cdots$ be the sets of
Theorem~\ref{thm:raise}. Let $D^{-}$ be the first $D_j$ that is at least $\sqrt x$, and let $D^{+}$
be the first $D_j$ that is at least $x/16$. For large $r$,
\[
  \sqrt x\le D^{-}<2\sqrt x,\qquad \frac x{16}\le D^{+}<\frac x8,\qquad D^{-}<D^{+}.
\]
Let $S^{-}\subseteq S^{+}$ be the corresponding sets, and put $J^{\pm}=G-(P\cup S^{\pm})$. Then
$J^{-}$ has girth greater than $D^{-}$, $J^{+}$ has girth greater than $D^{+}$, and
$J^{+}\subseteq J^{-}$. By Theorem~\ref{thm:raise}(ii),
\begin{equation}\label{eq:twosizes}
  |S^{-}|\;\le\;2^{4\sqrt x+K_1\log^{2}(2\sqrt x+2)}=r^{o(1)},\qquad
  |S^{+}|\;\le\;2^{x/4+K_1\log^{2}(x)}=r^{1/4+o(1)}.
\end{equation}
The small scale $D^{-}$ will make the final pieces clean. The large scale $D^{+}$ gives a long chain.

\paragraph{Cycles at the large scale.}
We apply Section~\ref{sec:chain} with $S=S^{+}$, $J=J^{+}$, $D=D^{+}$ and
$s=\lfloor(D^{+}-1)/2\rfloor\ge x/32-2$. The hypotheses of Lemma~\ref{lem:linkshalf} hold for large
$r$: $8\le D^{+}\le L$, and $2(t+2\lfloor D^{+}/8\rfloor)\le D^{+}$ because $t\le z\le D^{+}/4$. So
every two disjoint cycles of $J^{+}$ of length at most $L$ have at most
\[
  b\;=\;1+\frac{2L}{t+1}+2^{t+1}+578\Bigl(\frac L{D^{+}}\Bigr)^{2}
  \;\le\;1+\frac{2Ax}{z-2\log z}+\frac{2x}{z^{2}}+578\,(16A)^{2}\;\le\;\frac{3Ax}{z}
\]
links. Hence
\[
  \eta=2^{b}L^{3}2^{-s}\;\le\;r^{3A/z}\,(Ax)^{3}\,4\,r^{-1/32}=r^{-1/32+o(1)}=o(1),
  \qquad L^{2}2^{-s}=o(1).
\]

The graph $J^{+}$ is simple, with maximum degree three and girth greater than $D^{+}$. It has
$v=2r-2-|P|-|S^{+}|=2r-o(r)$ vertices, and its deficit is $d_J\le3(|P|+|S^{+}|)=r^{1/4+o(1)}$. Let
$\mu$ be the total weight of its cycles of length at most $L$. The proof of
Lemma~\ref{lem:weight} applies to $J^{+}$ with these parameters. It gives
$\mu\ge(1-o(1))\,T$, and for even $\ell$ with $6x\le\ell\le L$
\[
  2^{-\ell}\tr B^{\ell}\;\ge\;1-\frac{L\,d_J}{2v}-3v\,r^{-3}=1-o(1).
\]
The sum of $1/(2\ell)$ over the even $\ell$ in $[6x,L]$ is $(\frac14-o(1))\ln(A/6)=42-o(1)$. So
$\mu>20$ for large $r$.

\paragraph{The chain.}
Fix $p\in P$. By Lemma~\ref{lem:chain} there are a cycle $C$ of $J^{+}$ of length at most $L$, a
vertex $u$, and pairwise disjoint cycles $C_1,\dots,C_T$ of $J^{+}$ of length at most $L$ with
\[
  T\;\ge\;\frac{8\cdot2^{s}}{L^{3}}\;\ge\;\frac{2}{A^{3}}\cdot\frac{r^{1/32}}{x^{3}},
\]
such that each $C_i$ separates $u$ from $p$ in $G-V(C)$. Let $H$ be the component of $G-V(C)$
that contains $u$ and $p$. Each $V(C_i)$ is connected and separates $u$ from $p$, so it lies in $H$.
By Lemma~\ref{lem:nested} we may number the cycles so that $C_1<C_2<\dots<C_T$. Let $U_i$ be the
vertex set of the component of $u$ in $H-V(C_i)$.

\paragraph{Blocks.}
Put $w=\lceil128A^{2}x\rceil+2$. For $l=0,1,\dots$ with $(l+1)w\le T-1$ let
\[
  X_l\;=\;U_{(l+1)w+1}\sm\bigl(U_{lw+1}\cup V(C_{lw+1})\bigr).
\]
There are at least $\lfloor(T-1)/w\rfloor\ge r^{1/32}/x^{5}$ such sets, for large $r$. By
Lemma~\ref{lem:blocks} they are pairwise disjoint. Each contains the $w-1$ cycles between its two
end cycles. These cycles are pairwise disjoint and each has more than $D^{+}$ vertices, so
\[
  |X_l|\;\ge\;(w-1)(D^{+}+1)\;\ge\;128A^{2}x\cdot\frac x{16}=8A^{2}x^{2}\;\ge\;8L^{2}.
\]
Every edge of $H$ that leaves $X_l$ ends on one of its two end cycles. An end cycle has at most $L$
vertices, and each of them has one edge end outside the cycle. So at most $2L$ edges of $H$ leave
$X_l$.

\paragraph{Clean blocks.}
Discard a set $X_l$ if it contains a vertex with a neighbour on $C$, or a vertex of $S^{-}$. At most
$L$ vertices outside $C$ have a neighbour on $C$, so this discards at most
$L+|S^{-}|=r^{o(1)}$ sets. For large $r$ at least two sets remain. Let $X$ be one of them.
\begin{itemize}
\item $X$ avoids $P$. The set $P$ is connected and disjoint from the cycles, and it contains $p$,
which is not in $U_i$ for any $i$. A path in $G[P]$ from $p$ to a vertex of $U_i$ would have to pass
through $C_i$ or through $C$.
\item No edge of $G$ joins $X$ to $V(C)$, and $H$ is a component of $G-V(C)$. So the edges of $G$
that leave $X$ are edges of $H$, and $|\delta X|\le2L$.
\item $X$ avoids $P\cup S^{-}$, so $G[X]$ is a subgraph of $J^{-}$, and its girth is greater than
$D^{-}$.
\end{itemize}

\paragraph{The contradiction.}
The two remaining sets satisfy the hypotheses of Proposition~\ref{prop:blocks}(b) with
$\Lambda=L$ and $D=D^{-}$: they are disjoint, they avoid $P$, each has at least $8L^{2}$ vertices
and at most $2L$ boundary edges, and $D^{-}\ge\sqrt x\ge C_0\log(L+2)$ for large $r$. So
$q(G,P)\le\frac9{25}$. This contradicts $q(G,P)>\half$. Hence there is $R$ such that
$q(G,P)\le\half$ whenever $r\ge R$. \qed

\begin{remark}\label{rem:restore}
The set $X$ is defined in a component $H$ of $G-V(C)$, and not in $J^{+}$. So it contains the vertices of
$S^{+}\sm S^{-}$ that lie between its two end cycles. This is the reason for the two scales. The
chain needs the girth of $J^{+}$, and the pieces need only the girth of $J^{-}$, which has far fewer
deleted vertices than the chain has cycles.
\end{remark}

\begin{remark}\label{rem:polynomial}
The same chain proves a quantitative form of Theorem~\ref{thm:forest} for graphs with short cycles.
Suppose that $q(\Gamma,P)>\half+\varepsilon$ in the setting of Section~\ref{sec:forest}. Delete
short cycles at two scales with Lemma~\ref{lem:smallcuts}, find the chain, and count the weight of
the cycles inside the clean blocks. This weight is at most $(1+o(1))/(4\varepsilon)$ by the second
moment, and it is at least a constant times the number of clean blocks. The result is
$q(\Gamma,P)\le\half+r^{-\kappa}$ for every fixed $\kappa<\frac13$, when $|P|\le r/\log^{2}r$ and
$\Gamma[P]$ has at most $\gamma\log r$ components, for a constant $\gamma=\gamma(\kappa)>0$. Both
models checked this proof; it is in the repository as a research note~\cite{S4} and it is not
printed here. With an exponential form of Lemma~\ref{lem:smallcuts}, Astra's report~\cite{S4}
improves the error term to $O((\log r)^{11/2}/\sqrt r)$. Fable~5.2 has not checked that last
refinement line by line.
\end{remark}
